% Generated draft: requires mathematical and source review.
\documentclass[11pt]{article}
\usepackage[margin=25mm]{geometry}
\usepackage{fontspec}
\setmainfont{Latin Modern Roman}
\usepackage{amsmath,amssymb,mathtools,unicode-math}
\setmathfont{Latin Modern Math}
\usepackage{enumitem,needspace,xurl,hyperref}
\setlength{\parindent}{0pt}
\setlength{\parskip}{5pt}
\setlength{\emergencystretch}{3em}
\newcommand{\N}{\mathbb{N}}
\newcommand{\Z}{\mathbb{Z}}
\newcommand{\Q}{\mathbb{Q}}
\newcommand{\R}{\mathbb{R}}
\newcommand{\C}{\mathbb{C}}
\newcommand{\F}{\mathbb{F}}
\newcommand{\A}{\mathbb{A}}
\newcommand{\PP}{\mathbb{P}}
\newcommand{\E}{\mathbb{E}}
\newcommand{\eps}{\varepsilon}
\newcommand{\dd}{\,\mathrm{d}}
\begin{document}
\section*{UnsolvedMath ID: 20001048}
\textbf{Source problem code:} catalogue entry
\subsection{Definitions and mathematical statement}
The exact assigned source text is only ``Conjecture 6.8. Under the hypotheses of''. The verification field says that the canonical text is split across two consecutive extraction records, but neither the referenced hypotheses nor the continuation is included here. Accordingly, the complete original conjecture cannot be stated from this record alone.

The supplied research aid calls its target a directed overlapping-list rainbow-cycle conjecture. To express its reported structural content, use the following explicit \emph{conditional model}, not a verified reconstruction of Conjecture 6.8. Let \(V\) have size \(n\geq2\), let \(I=\{1,\ldots,r\}\) be a finite label set, and let \(E_i\subseteq\{(u,v)\in V^2:u\ne v\}\) be the arcs carrying label \(i\). Different labels may occur on the same ordered pair; antiparallel arcs are allowed in this model. Put
\[
 E=\bigcup_{i\in I}E_i,\qquad L(e)=\{i\in I:e\in E_i\},\qquad
 d_i^+(v)=|\{w:(v,w)\in E_i\}|.
\]
A rainbow directed cycle is a directed cycle with distinct vertices, together with an injective choice of one label from \(L(e)\) for each of its arcs. A possible threshold formulation suggested by the summary is: if \(d_i^+(v)\geq d_i\) for every \(v,i\), where the \(d_i\) are nonnegative integers and \(\sum_i d_i\geq n\), then a rainbow directed cycle exists. The supplied fragment does not establish that these are exactly the original hypotheses.

In this terminology the aid reports the following counterexample normalization: after deleting label incidences and repeatedly restricting to sink strongly connected components, one obtains a counterexample on \(V'\), of size \(n'\), whose union digraph is strongly connected and for which
\[
 d_i^+(v)=d_i\quad(v\in V',\ i\in I),\qquad \sum_i d_i=n'.
\]
Here zero layers may be omitted, and all quantities refer to the normalized instance. Define its overlap core by
\[
 H=(V',\{e:|L(e)|\geq2\}).
\]
The reported conclusions are \(\delta^+(H)\geq1\), and, for every directed cycle \(C\) of \(H\), there is \(S\subseteq E(C)\) with
\[
 |S|\geq3,\qquad \left|\bigcup_{e\in S}L(e)\right|<|S|.
\]
Thus the obstruction is failure of a system of distinct representatives on that cycle. The aid further reports that \(H\) has no directed \(2\)-cycle and that, on each directed triangle \(e_1,e_2,e_3\) of \(H\),
\[
 L(e_1)=L(e_2)=L(e_3)=\{a,b\}\quad\text{for some distinct }a,b\in I.
\]
These are reported necessary features of normalized counterexamples, not a proof that no counterexample exists.

The reported sharpness example can be expressed in the same conditional model: on a directed \(n\)-cycle, give each arc every label in \(\{1,\ldots,n-1\}\). Each layer has outdegree one, total layer degree is \(n-1\), and the sole directed cycle has too many arcs for a rainbow labeling. This illustrates the claimed threshold, subject to confirming the missing original hypotheses.
\subsection{Short English statement}
The assigned source contains only the beginning of a conjecture. The accompanying aid describes a reduction of possible rainbow-cycle counterexamples to tight strongly connected instances with a constrained overlap core; the original hypotheses still need recovery.
\subsection{Sources}
\begin{itemize}
\item Local registry: permanent UnsolvedMath ID \(20001048\), exact public code \texttt{catalogue entry}; \texttt{aim-combinatorics-notes.json}, \texttt{notes[172]}.
\item AIM workshop ``The Caccetta-Haggkvist conjecture'', source item 6.8, canonical tag conjecture: \url{https://aimath.org/WWN/caccetta/caccetta.pdf}. The supplied extracted text is ``Conjecture 6.8. Under the hypotheses of''. The literature field and remarks are empty. No proposer or year is supplied.
\item The supplied statement-verification note reports a split across consecutive extraction records. The adjacent continuation is not supplied in this record.
\item Research attempt 1 reports a \texttt{partial\_result} and structural reduction, with status at the run \texttt{open}, rigor label \texttt{proved}, novelty confidence low, classification \texttt{PARTIAL-PROGRESS}, and suggested difficulty L4. It provides no separate authoritative citation for the recovered formulation or reduction.
\end{itemize}
\end{document}
